\chapter{Conclusão}
\label{ch:conclusao}

Este capítulo retoma as atividades propostas na seção 1.6 para avaliar sua execução, apresenta os resultados obtidos ao final do trabalho e descreve o planejamento para a fase seguinte da monografia.

\section{Atividades Propostas}

\begin{table}[h]
    \centering
	\caption{Atividades Concluídas}
    \begin{tabular}{|p{9cm}|p{1.8cm}|}
        \hline
        \textbf{Atividade} & \textbf{Situação} \\
        \hline
        Contextualização sobre Engenharia de Software Experimental & Concluída  \\
        \hline
        Definição do GQM & Concluída \\
        \hline
        Elaboração do Protocolo de Revisão da Literatura & Concluída \\
        \hline
        Execução da String de Busca e Seleção dos Artigos & Concluída \\
        \hline
        Leitura do Material Selecionado & Concluída \\
        \hline
        Definição da Proposta de Solução & Concluída \\
        \hline
        Escrita da Monografia & Concluída \\
        \hline
        Revisão da Monografia & Concluída \\
        \hline
        Apresentação da Monografia & Agendada \\
        \hline
    \end{tabular} \\[0.5em]
   	Fonte: Autor
	\label{tab:atividades-concluidas}
\end{table}

\section{Resultados Obtidos}

A revisão estruturada da literatura permitiu consolidar o estado da arte sobre o uso de especificações BDD em conjunto com técnicas de Inteligência Artificial para apoio à automação de testes de software. A partir da análise dos 22 estudos selecionados, foi possível identificar evidências de que a forma de escrita, a estruturação e o nível de detalhamento dos artefatos textuais influenciam a capacidade de modelos baseados em IA de gerar cenários e códigos de teste com maior aderência ao comportamento esperado.

Os estudos analisados também evidenciaram que, embora existam avanços relevantes na automação de testes com IA, ainda persistem limitações relacionadas à ambiguidade da linguagem natural, ao contexto insuficiente das especificações e à dificuldade de traduzir requisitos escritos em artefatos executáveis de forma consistente. Em contrapartida, a literatura também mostrou que cenários mais estruturados e contextualizados tendem a favorecer a geração de artefatos mais completos e com menos falhas, especialmente quando apoiados por estratégias de prompting e por mecanismos adicionais de contextualização.

Com base nessa revisão, o principal resultado desta primeira etapa foi a consolidação do protocolo metodológico do estudo observacional retrospectivo a ser executado no contexto do CAPJu. Esse protocolo passou a incluir a definição do caso, da questão específica do estudo, das fontes de dados, das unidades de análise e das métricas utilizadas para avaliar a eficácia e a qualidade do código de testes automatizados gerado por IA a partir de especificações BDD.

Dessa forma, esta etapa da monografia forneceu não apenas a fundamentação teórica do problema, mas também a estrutura metodológica necessária para a execução empírica do estudo na segunda etapa do trabalho. Espera-se, com isso, produzir evidências sobre como a forma de escrita das especificações BDD se relaciona com os resultados obtidos na geração automatizada de testes, contribuindo para a compreensão do uso de IA nesse contexto.

\section{Cronograma da Segunda Etapa da Monografia}

O cronograma que pode ser visto na Figura de atividades \ref{fig:cronograma_2} está organizado para o período de março a julho, iniciando com a evolução do trabalho a partir das observações e a preparação dos artefatos nos meses de março e abril. Em seguida, entre abril e maio, é realizada a execução do estudo observacional retrospectivo, com posterior análise e consolidação dos resultados ao longo de maio e junho. A escrita da monografia acontece principalmente em junho, seguida pela fase de revisão entre junho e julho, finalizando com a apresentação da monografia no mês de julho, concluindo assim todas as etapas do trabalho.
